\input{preamble}
\usepackage{graphicx}
\usepackage{url}

\begin{document}

\title{Nakajima's representation of the Heisenberg algebra}
\author{Daniel González Casanova Azuela}

\maketitle

\phantomsection
\label{section-phantom}

\section*{Abstract}
\label{section-abs}

\noindent
The main theorem in \cite{Nakajima-Heis}
shows there is an action of the infinite-dimensional
Heisenberg algebra
on the homology groups of the Hilbert
scheme of points on a projective surface.
We explain how this representation is constructed
and give an idea of the proof of the theorem.



\tableofcontents

\section{Introduction}
\label{section-intro}

\noindent
Motivated by earlier work on moduli spaces of 
torsion-free sheaves
(see \cite{Nakajima-ALE}),
Nakajima constructed in \cite{Nakajima-Heis}
a representation the infinite-dimensional Heisenberg
algebra on the homology groups of the Hilbert schemes
on a nonsingular projective surface.

In this work we outline the basic ideas in Nakajima's construction
and explain the first part of the proof of the main theorem.
Most of this work is based in \cite{Nakajima-lectures},
with the exception of the first section
where I try to identify the precise version of the Heisenberg
algebra that is being represented.
In Section \ref{section-hilbert-scheme},
I explain the construction of the Hilbert scheme,
providing slightly more details than what is really essential 
to understand the main theorem,
but which help to get a better idea on what the Hilbert scheme is
and how to work with it.

This work is strongly influenced by \cite{videos}, 
a series of seminars aimed at understanding \cite{Nakajima-lectures}
documented in videos.
This was of great help
specially in Section \ref{section-bm}, since
the properties of Borel-Moore homology are a rather technical
but essential part of the construction.
The remaining sections, containing the main theorem
and part of its proof,
also pick up ideas from these videos,
but are mostly based in \cite{Nakajima-lectures}.


\section{Heisenberg algebra}
\label{section-heisenberg}

\noindent
In this section we review four different variants
of the Heisenberg algebra and discuss their characters.
We conclude by showing how the character is used
to prove that Nakajima's representation is
of highest weight.

\medskip\noindent
The most basic definition of the Heisenberg Lie algebra is
$\hat{\mathfrak{a}}=
\bigoplus_{n \in \mathbb{Z}} \mathbb{C}h_n \oplus\mathbb{C}K$
with bracket given by 
\begin{equation}
\label{equation-bracket-basic}
[h_m,h_n]=m\delta_{m,-n}K,\qquad 
[K,\hat{\mathfrak{a}}]=0.
\end{equation}

\noindent
In class we described its highest weight representation
$H=\mathbb{C}[x_1,x_2,\ldots]$ with
\begin{align*}
h_n&\mapsto \begin{cases}
n \frac{\partial }{\partial x_n}\qquad &\text{if }n>0\\
x_{-n}\qquad &\text{if }n<0\\
0\qquad &\text{if }n=0
\end{cases},
\qquad \qquad  K\mapsto \text{Id}.
\end{align*}

\noindent
We introduced a grading in $H=\bigoplus_{\Delta=0}^\infty H_\Delta$ 
given by the {\it energy} of monomials, defined by
$$
\Delta(x_1^{k_1}x_2^{k_2}\ldots x_s^{k_s})=\sum_j k_j,
$$
which gives the character formula
\begin{equation}
\label{equation-character-heisenberg}
\sum_{\Delta=0}^\infty q^\Delta \dim H_\Delta
=\prod_{m=1}^\infty\frac{1}{1-q^m}.
\end{equation}

\noindent
In order to generalize this construction
it's worth going over why this formula holds.
Recall that the \textit{Verma module} of
weight $\mu=0$
of $\hat{\mathfrak{a}}$ is defined by
$H=
\mathcal{U}(\hat{\mathfrak{a}})
\otimes_{\mathcal{U}(\hat{\mathfrak{a}}_+)}
\mathbb{C}|0\rangle$.
Further, by the PBW theorem, this module
is isomorphic to $\mathcal{U}(\mathfrak{n}_-)$,
the universal enveloping algebra
of the negative roots in the triangular decomposition
(as in the finite-dimensional case, see \cite[Lecture 24]{KLAL} and
\cite[Video 8]{videos}).
Since the subalgebra of negative roots is abelian,
the universal enveloping algebra becomes the symmetric algebra
and therefore we have
\begin{equation}
\label{equation-symmetric-rank-1}
\begin{aligned}
H&\cong\text{Sym}\left(\bigoplus_{m>0}\mathbb{C}at^{-m}\right)\\
&\cong\bigotimes_{m>0}\text{Sym}(\mathbb{C}at^{-m})\\
&\cong \bigotimes_{m>0}\mathbb{C}[at^{-m}],
\end{aligned}
\end{equation}

\noindent
since the symmetric power of a vector space
is the polynomial ring in its generators,
and the symmetric power of a direct sum
is the tensor product of the symmetric powers of the factors.
This matches our previous formula
$H=\mathbb{C}[x_1,x_2,\ldots]$ by letting $x_m=at^{-m}$.

Putting a dummy variably $q$ to every 1-dimensional
factor and multiplying throughout, we obtain
\begin{align*}
m=1\qquad &(1+q+q^2+q^3+\ldots)\times\\
m=2\qquad &(1+q^2+q^4+q^6+\ldots)\times\\
m=3\qquad &(1+q^3+q^6+q^9+\ldots)\times\ldots,
\end{align*}

\noindent
so that indeed the character is given by 
Equation \ref{equation-character-heisenberg}.

\medskip\noindent
Another version of this construction
we discussed in lectures
(see \cite[Section 3.1]{Jethro-vertex-algebras})
starts with any finite-dimensional vector space $\mathfrak{h}$
equipped with a symmetric bilinear form $(\cdot,\cdot)$.
Define
$$
\hat{\mathfrak{h}}=\mathfrak{h}\otimes\mathbb{C}[[t^{\pm 1}]]\oplus\mathbb{C}K
$$
with bracket
\begin{equation}
\label{equation-bracket-affine}
[at^m,bt^n]=m\delta_{m,-n}(a,b)K,\qquad [K,\hat{\mathfrak{h}}]=0,
\end{equation}

\noindent
which is just the affine Lie algebra of $\mathfrak{h}$ when
regarded as a commutative Lie algebra.

A similar formula to Equation \ref{equation-symmetric-rank-1} is
\begin{equation}
\label{equation-symmetric}
\begin{aligned}
H&\cong\text{Sym}\left(\bigoplus_{m>0}\bigoplus_{a \in B}
\mathbb{C}at^{-n}\right)\\
&\cong\bigotimes_{m>0}\bigotimes_{a \in B}\text{Sym}(\mathbb{C}at^{-m})\\
&\cong\bigotimes_{m>0}\bigotimes_{a \in B}\mathbb{C}[at^{-m}]
\end{aligned}
\end{equation}
where $B$ is a basis of the underlying vector space $\mathfrak{h}$.
If $\mathfrak{h}$ has dimension $r$, this has character
\begin{equation}
\label{equation-character-rank-r}
\prod_{m=0}^\infty\frac{1}{(1-q^m)^r}.
\end{equation}

\medskip\noindent
Taking this construction one step further,
let's change $\mathfrak{h}$ for a finite-dimensional
super vector space $\mathfrak{h}=\mathfrak{h}_0\oplus \mathfrak{h}_1$
with basis $B^{\text{even}}\amalg B^{\text{odd}}$.
If we equip $\mathfrak{h}$ with a graded commutative
pairing $\left<\cdot,\cdot\right>$, we can
apply the same construction as above
to obtain the algebra
\begin{equation}
\label{equation-super-heisenberg}
\hat{\mathfrak{h}}
=(\mathfrak{h}_0[[t^{\pm 1}]]\oplus\mathfrak{h}_1[[t^{\pm 1}]])
\oplus\mathbb{C}K
\end{equation}

\noindent
with Lie bracket
\begin{equation}
\label{equation-super-bracket}
[at^m,bt^n]=\left<a,b\right>m\delta_{m,-n}K,\qquad 
[K,\hat{\mathfrak{h}}]=0.
\end{equation}

\noindent
One can check that this is indeed a super Lie bracket
(which essentially holds by the pairing being graded commutative
and $K$ being central),
i.e. $\hat{\mathfrak{h}}$ is a Lie superalgebra.
Its highest weight module is given by
\begin{align*}
H&\cong \text{Sym}\left(\bigoplus_{m>0}\bigoplus_{a \in B}
\mathbb{C}at^{-m}\right)\\
&\cong \bigotimes_{m>0}\text{Sym}\left(\bigoplus_{a \in B^{\text{even}}}
\mathbb{C}at^{-m}\right)
\otimes\bigotimes_{m>0}
\Lambda\left(\bigoplus_{a \in B^{\text{odd}}}
\mathbb{C}at^{-m}\right).
\end{align*}

\noindent
since the universal enveloping algebra of the odd
part is the exterior power (see \cite[Video 8]{videos}).

Now recall that while the symmetric power of a vector
space yields the polynomial algebra on the basis of the vector space,
the exterior power yields the algebra generated
by the wedge products of the basic vectors.
In the case of a 1-dimensional vector space generated
by the element $at^{-m}$, we get
$$
\Lambda(\mathbb{C}at^{-m})=\mathbb{C} \oplus \mathbb{C}at^{-m}
$$
and thus the character contribution in this case is simply
$
1+q^m.
$
Multiplying throughout we obtain the character
$$
\prod_{m=0}^\infty \frac{(1+q^m)^{\dim \mathfrak{h}^{\text{odd}}}}{
(1-q^m)^{\dim \mathfrak{h}^{\text{even}}}}.
$$

\medskip\noindent
Let's introduce one last variant of this construction.
This time we start from a graded vector space
$\mathfrak{h}=\bigoplus_{m \in \mathbb{Z}}\mathfrak{h}_m$,
which is automatically also a supervector space
by the $\mathbb{Z}_2$-grading
$\mathfrak{h}^{\text{even}}=\bigoplus_{i \in \mathbb{Z}}\mathfrak{h}_{2i}$ 
and $\mathfrak{h}^{\text{odd}}=\bigoplus_{i \in \mathbb{Z}}\mathfrak{h}_{2i+1}$.
Suppose $\mathfrak{h}_i$ is finite-dimensional with basis $B_i$,
and that $\mathfrak{h}$ is equipped with a graded commutative
pairing $\left<\cdot,\cdot\right>$.
The definition of $\hat{\mathfrak{h}}$ as an algebra
is given exactly as in Equation \ref{equation-super-heisenberg}
with bracket given by Equation \ref{equation-super-bracket}.

Now we obtain
\begin{align*}
H&\cong \text{Sym}\left(\bigoplus_{m>0}\bigoplus
_{\substack{a \in B_i \\ i \in \mathbb{Z}}}
\mathbb{C}at^{-m}\right)\\
&\cong \bigotimes_{m>0}\text{Sym}\left(\bigoplus
_{\substack{a \in B_{2i} \\ i \in \mathbb{Z}}}
\mathbb{C}at^{-m}\right)
\otimes\bigotimes_{m>0}
\Lambda\left(\bigoplus_{
\substack{a \in B_{2i+1} \\ i \in \mathbb{Z}}}
\mathbb{C}at^{-m}\right).
\end{align*}

\noindent
This gives the character
$$
\sum_{\substack{\Delta\in \mathbb{Z}_{\geq 0}  \\ i \in \mathbb{Z}}}
\dim \mathfrak{h}_{i,\Delta}t^iq^\Delta
=\prod_{m>0}\prod_j(1-(-1)^jt^jq^m)^{-(-1)^j\dim \mathfrak{h}_j}.
$$

\noindent
Finally fix $\mathfrak{h}=H_*(X)$, the singular homology ring
of a complex projective surface $X$.
Then the character of the corresponding Heisenberg algebra is
\begin{equation}
\label{equation-character}
\prod_{m=1}^\infty \frac{
(1+t^{2m-1}1^m)^{b_1(X)}
(1+t^{2m+1}q^m)^{b_3(X)}}
{(1-t^{2m-2}q^m)^{b_0(X)}
(1-t^{2m}q^m)^{b_2(X)}
(1-t^{2m+2}q^m)^{b_4(X)}}.
\end{equation}

\noindent
It was proved by Göttsche using the Weil conjectures in \cite{Gottsche1990} 
(and in \cite[Section 6.2]{Nakajima-lectures} by other methods)
that Equation \ref{equation-character}
in fact equals
$$
\sum_{n=0}^\infty q^nP_t(X^{[n]}),
$$
where the {\it Poincaré polynomials} $P_t(Z)$ are defined by
$P_t(Z)=\sum_{m\geq 0}b_m(Z)t^m$, $b_m(Z)=\text{rank}H_m(Z)$ for
any variety $Z$.

The representation
of the Heisenberg algebra we discuss in this work corresponds
precisely to this version of the Heisenberg algebra.
Göttsche's formula shows that it is a highest weight representation,
as explained in \cite[Corollary 8.16]{Nakajima-lectures}.




\section{Hilbert scheme of points}
\label{section-hilbert-scheme}

\noindent
In this section we define 
of the Hilbert scheme of points
on a smooth projective complex surface $X$.
The essential property of the Hilbert scheme
is that its points (over $\mathbb{C}$) 
are 0-dimensional closed subschemes of $X$
of length $n$.

\medskip\noindent
Consider a contravariant functor 
$$
\mathcal{H}ilb_X:(\text{Schemes})^{\text{op}}\to (\text{Sets})
$$
which assigns to any scheme $U$ the set of flat families
of closed subschemes in $X$ parametrized by $U$.
Any such family is by definition a flat morphism
$$
\xymatrixcolsep{.1em}
\xymatrix{
Z \ar[d]^\pi&\subseteq X \times U\\
U.
}
$$
We define $\mathcal{H}ilb$ to map a morphism of schemes
$U \to U'$ to the map of sets (of families over $U$ and $U'$) given
by pullback. That is, we map a flat family $Z'\to U'$ 
to $Z'\times_{U'}U \to U$, which remains flat.

The Hilbert polynomial at a fibre $Z_u$ is defined by
$P_u(m)=\chi(\mathcal{O}_{Z_u}\otimes \mathcal{O}_X(m))$.
By flatness of $\pi$ the Hilbert polynomial is constant along $U$
(\cite[Part III, Theorem 9.9]{Hartshorne}).
For a fixed polynomial $P$
we define $\mathcal{H}ilb^P_X$ as the functor which maps $U$ 
to the set of families over $U$ with Hilbert polynomial $P$.
A theorem by Grothendieck says that $\mathcal{H}ilb^P_X$ is representable
by a projective scheme $\text{Hilb}^P_X$.
This means that there is a natural transformation
$\Hom(-,\text{Hilb}^P_X) \simeq \mathcal{H}ilb^P_X(-)$;
that is, families over any scheme $U$ are in bijection with
morphisms $U \to \text{Hilb}^P_X$.

As with any representable functor,
we obtain a distinguished element associated to the identity morphism
of the object representing the functor.
In our case, this is a
``universal'' family $\mathcal{Z} \to \text{Hilb}^P_X$
with the property that any other family $Z \to U$ is obtained as a pullback
of $\mathcal{Z}$.
More exactly, when regarding an arbitrary family $Z \to U$ as a morphism
$f \in \Hom(U,\text{Hilb}^P_X)$, we see that
$$
\xymatrix{
\text{id}_{\text{Hilb}^P_X}\ar@{|->}[d]
&\Hom(\text{Hilb}^P_X,\text{Hilb}^P_X)\ar[r]^-{\simeq}\ar[d]_{f^*}
&\mathcal{H}ilb^P_X(\text{Hilb}^P_X)\ar[d]^{U\times_{\text{Hilb}_X^P}-}
&\mathcal{Z}\ar@{|->}[d]\\
f&\Hom(U,\text{Hilb}^P_X)\ar[r]^-{\simeq}
 &\mathcal{H}ilb^P_X(U)& Z=U\times_{\text{Hilb}_X^P}\mathcal{Z}.
}
$$
commutes by naturality.
This shows that studying families of subschemes of $X$ parametrized by any
scheme $U$ can be done by studying the universal family
over the Hilbert scheme.

We can formally view subschemes of $X$
as points of $\operatorname{Hilb}_X^P$ as follows.
Recall that, as with any other scheme,
points $h$ of $\text{Hilb}_X^P$ correspond to morphisms
$\Spec(\kappa(h)) \to \text{Hilb}^P_X$, where $\kappa(h)$ is the residue field of the point
(see Stacks Project 
\href{https://stacks.math.columbia.edu/tag/01J9}{01J9}).
Then for every point $h$ of the Hilbert scheme
we have a unique subscheme
of $X$ defined as the fibre $X_h$
of the corresponding morphism,
$$
\xymatrix{
X_h\ar[r]\ar[d]\ar@{}[dr]|<{\lrcorner}&\mathcal{Z}\ar[d]
& \subset X \times \text{Hilb}^P_X\\
\Spec(\kappa(h))\ar[r]&\text{Hilb}^P_X.
}
$$
As a rather technical observation that will not be
addressed later, we note that 
to make sure that the fibres are $\mathbb{C}$-schemes,
we must restrict our analysis to $\mathbb{C}$-points
of the Hilbert scheme (i.e. when the residue field is $\mathbb{C}$).

\begin{definition}
\label{definition-hilbert-scheme-of-points}
Let $P$ be a constant polynomial given by $P(m)=n$ for all $m \in \mathbb{Z}$.
The {\it Hilbert scheme of $n$ points in $X$} is $X^{[n]}:=\text{Hilb}^P_X$.
\end{definition}

\noindent
This condition forces the dimension of the fibres of any family
over $X$ to be zero. Indeed, the leading term
of the Hilbert polynomial of a projective variety of degree $d$
and dimension $s$ is $(dm^s)/s!$, which can only be constant
for $s=0$ (see \cite[Exercise 1.13]{HarrMorr}).

Further, any 0-dimensional subscheme of $X$
is supported on finitely many points of $X$
(this is consequence of $X$ being Noetherian:
we couldn't have infinitely many points since they
would give an infinite chain of prime ideals).
Finally, the {\it length} of any such subscheme, defined by
 $\dim\Gamma(\mathcal{O}_Z)$
must be $\sum_{x \in \text{supp} Z}\dim\Gamma(\mathcal{O}_{Z_x})=n$.
(see Stacks Project 
\href{https://stacks.math.columbia.edu/tag/0AYT}{0AYT}).

\medskip\noindent
It is now easy to see that any subscheme of $X$ consisting of $n$ distinct
points is itself a point of $\text{Hilb}_X^P$.
Indeed, if $Z=\{x_1,\ldots,x_n\}\subset X$, its
structure sheaf is the direct sum of skyscraper sheaves
$$
\mathcal{O}_Z=\bigoplus_{i=1}^n (\text{skyscraper sheaf at $x_i$})
$$
so that $\mathcal{O}_Z \otimes \mathcal{O}_X(m)\cong\mathcal{O}_Z$
and thus the Hilbert polynomial of $Z$ is constant $n$.
The latter isomorphism can be verified via
the exact sequence associated to the maximal ideal 
$\mathfrak{m}_{x_i}$
of the stalk $\mathcal{O}_{X,x_i}$, namely
$$
\xymatrix{
0\ar[r]
&\mathfrak{m}_{x_i}\ar[r]
&\mathcal{O}_{X,x_i}\ar[r]
&\mathcal{O}_{X,x_i}/\mathfrak{m}_{x_i}\cong\mathbb{C}\ar[r]
&0,
}
$$
which after tensoring by 
$\mathcal{O}_{Z,x_i}
\cong\mathbb{C}\cong\mathcal{O}_{X,x_i}/\mathfrak{m}_{x_i}$ 
kills the first term.
Twisting by $m$ does not alter this construction.

\medskip\noindent
However, not every point in $X^{[n]}$ is given
by a set of $n$ distinct points,
nor by a set of points with multiplicites,
as the following construction suggests.
Define the {\it symmetric power} of $X$ as
$$
S^nX=\underbrace{X \times \ldots \times X}_{n\text{ copies}}/S_n
$$
where the symmetric group $S_n$ acts by
permutating copies. Thus, a point of $S^nX$ is
an $S_n$-orbit on the cartesian product $X \times \ldots \times X$,
i.e., an unordered $n$-tuple of points of $X$
with allowed repetition.
In algebraic geometry these are called ``algebraic cycles''
and are denoted as formal sums $\sum_{x \in X}n_x[x]$ 
where $\sum_{x \in X}n_x=n$, with $n_x \in \mathbb{N}$,
and $[x]$ are formal symbols.

By our discussion above on the length of 0-dimensional subschemes of $X$,
we have a map
\begin{align*}
\pi : X^{[n]} &\longrightarrow S^nX \\
 Z&\longmapsto \sum_{x \in \text{supp}Z}\dim
\Gamma(\mathcal{O}_{Z,x})[x]
\end{align*}
called the {\it Hilbert-Chow morphism}.
In some cases this is an isomorphism, but not always.

\begin{theorem}[Fogarty, \cite{Fogarty}]
\label{theorem-fogarty}
Suppose $X$ is nonsingular and $\dim X=2$,
then the following holds.
\begin{enumerate}
\item $X^{[n]}$ is nonsingular and has dimension $2n$.
\item The Hilbert-Chow morphism is a resolution of singularities.
\end{enumerate}
\end{theorem}




\section{Borel-Moore homology}
\label{section-bm}

\noindent
In this section we define Borel-Moore (BM) homology
and formulate its basic properties.
Most of this section is based on \cite[Video 9]{videos}.
For further details on the construction see \cite{Fulton-Young-Tableaux}.

\medskip\noindent
Let $M$ be a a smooth, connected,
oriented, finite-dimensional manifold over $\mathbb{R}$.
While the following results hold for any
topological space $X \subset M$
that is a proper deformation retract
of a closed neighbourhood of a manifold,
for us it's enough to think of $X \subset M$ as a
(possibly singular) variety over $\mathbb{C}$.

In this work we consider homology and cohomology
groups with coefficients in $\mathbb{C}$.
Recall that singular homology is defined
as the homology of the chain complex of
finite linear combinations of $i$-simplices in $X$.
\textit{Borel-Moore homology} $H_\bullet^{BM}(X)$ 
is the homology of the complex of locally finite
linear combinations of singular simplices in $X$,
where ``locally finite'' means that every
point $x \in X$ is in a neighbourhood
that intersects only finitely many 
simplices.

A striking 
difference between singular and BM homology
is that the latter is not homotopy invariant.
To see this recall that the homology groups
of  $\mathbb{R}^2$ are $H_0(\mathbb{R}^2)=\mathbb{C}$
and $H_i(\mathbb{R}^2)=0$ for $i \neq 0$.
However, $H_2^{BM}(\mathbb{R}^2)=\mathbb{C}$ 
since a tiling of triangles
is a locally finite chain of 2 simplices,
has no boundary (since the boundaries of
all triangles cancel each other),
and is clearly not the boundary of a 3-simplex;
thus this tiling is a nonzero element of $H_2^{BM}(\mathbb{R}^2)$.
Further, notice that $H_0^{BM}(\mathbb{R}^2)=0$
since a ray of 1-simplices
is a locally finite chain
whose boundary is a point;
thus a point is a boundary and 
has trivial homology class.

\medskip\noindent
Here are three important properties of BM homology:
\begin{enumerate}
\item If $X$ is compact, $H_i^{BM}(X)\cong H_i(X)$.
\item If $X$ is not compact,
$H_i^{BM}(X)\cong H_i(\hat{X},\infty)$ 
where $\hat{X}$ is the one-point
compactification of $X$ and
$H_i(\hat{X},\infty)$ is the
relative homology group.
\item 
\label{item-isomorphism-relative-cohomology}
If the dimension of the manifold
where $X$ is embedded is $m$,
there is a canonical isomorphism
$H_i^{BM}(X) \cong H^{m-i}(M,M\setminus X)$.
In other words, there is a perfect
pairing
$$
H_i^{BM}(X) \times H_{m-i}(M,M\setminus X)\to \mathbb{C}
$$
given by ``intersection''
(where the quotes account for some
transversal intersection condition
that we would require before
computing the intersection of
two simplicial chains).
\end{enumerate}

Consider the previous example of
the tiling of $\mathbb{R}^2$
by triangles and put 
$X=M=\mathbb{R}^2$.
Take any point as a 0-chain in singular homology
and the triangle tiling as 2-chain in BM homology.
Since the point lies in exactly one triangle
(after suitable perturbation)
their intersection is 1.
Notice this wouldn't work with
singular homology: we need an infinite 2-chain
to make sure that the whole plane
is tiled by triangles since
otherwise we cannot guarantee that the point
is contained in some triangle.

\medskip\noindent
Now we formulate the key properties
of BM homology needed to define the operators
in Nakajima's representation of the Heisenberg algebra.

\begin{enumerate}
\item (Pushforward.)
\label{item-pushforward}
If $f:X \to Y$ is proper
(i.e. inverse image of any compact set is compact)
then there is a pushforward map $f_* :H_i^{BM}(X)\to H_i^{BM}(Y)$.
This is functorial in proper maps,
i.e. preserves compositions.
Roughly, this works because, by properness, the inverse image of
a point in $Y$ will be mapped to a compact set of $X$,
which can intersect only finitely many simplices of a given chain,
ensuring local finiteness of the image under $f$ of the chain.

\item 
\label{item-pullback}
(Pullback.) If $f:X \to Y$ 
is a locally trivial fibration with
fibre a manifold of dimension $m$ (over $\mathbb{R}$),
then there is a pullback map
$f^* :H_i^{BM}(Y)\to H_{i+m}^{BM}(X)$.

\item  
\label{item-intersection}
(Intersection pairing.) If $X$ and $Y$ 
are both ``nicely embedded'' in the same
manifold $M$, there is a bilinear graded-commutative map
$$
\cap_M:H_i^{BM}(X)\times H_j^{BM}(Y)
\to H_{i+j-m}^{BM}(X\cap Y).
$$
Constructing this product can be done
via the isomorphism with relative singular cohomology
(item (\ref{item-isomorphism-relative-cohomology})
in the previous list),
and then using usual cap product.
Again, the term ``nicely embedded'' accounts
for a transversality condition that we will not
specify.

\item (Fundamental classes.)
\label{item-fundamental-class}
To any variety $V$ we can associate a fundamental class $[V]$.
More precisely,
\begin{itemize}
\item If $V$ is a smooth irreducible variety
of dimension $d$ over $\mathbb{C}$,
its {\it fundamental class} $[V]$ is the canonical generator of 
$H_{2d}^{BM}(V)\cong H^0(V)\cong \mathbb{C}$,
where the first isomorphism
is given by the isomorphism of BM homology and
relative singular cohomology taking $V = M$ since it is smooth.

\item If $V$ is a (possibly singular) irreducible variety of dimension $d$
over $\mathbb{C}$, its fundamental class $[V] \in B_{2d}^{BM}(V)$
is the unique class such that
$u^* [V]=[V_{\text{reg}}]$.
This is well-defined since
the pullback in this degree turns
out to be injective (by further properties of BM homology not discussed here)
and thus this class is unique.

\item For $V$ not necessarily irreducible
of dimension $d$, then $H_{2d}^{BM}(V)$ 
has basis $[V_1],\ldots,[V_k]$
where $V_i$ are the $d$-dimensional 
irreducible components of $V$.
In this case we define the fundamental class
as the sum of the fundamental classes of those
components, i.e. $[V]=[V_1]+\ldots+[V_k]$.
\end{itemize}
\end{enumerate}

\noindent
As a final remark notice that it is
as special feature of BM homology
that we can define fundamental classes for
noncompact manifolds, as opposed to 
singular homology.

\section{Correspondences}
\label{section-correspondences}

\noindent
In this section we introduce the main tool
for constructing the operators in Nakajima's representation
of the Heisenberg algebra.

\medskip\noindent
Let $X_1,X_2$ be smooth irreducible varieties
of dimensions $d_1,d_2$.
Let $Z$ be a $d$-dimensional 
closed subvariety of $X_1\times X_2$ 
such that the projection onto the first factor
is proper.

The {\it correspondence} associated to $Z \subset X_1\times X_2$
is the map
\begin{equation}
\label{equation-correspondence}
\begin{aligned}
c_{[Z]}:H_i^{BM}(X_2)&\longrightarrow H_{i+2d-2d_2}^{BM}(X_1)\\
c&\longmapsto p_{1,*}(p_2^*c\cap [Z])
\end{aligned}
\end{equation}

\noindent
where $p_1$ and $p_2$ are the projections in the following diagram:
$$
\xymatrix{
&Z\ar[ld]_{p_1}\ar[rd]^{p_2} \subset X_1\times X_2\\
X_1&&X_2.
}
$$
\noindent
That is, the correspondence is an operator
given by pulling back a class by $p_2$,
intersecting with the fundamental class of $[Z]$,
and pushing forward by $p_1$.
This map is well-defined according
to the key properties of BM homology stated in the previous section;
the degree of the target homology group is computed according
to such properties.

\begin{remark}
\label{remark-notes}
\begin{enumerate}
\item 
\label{item-any-class}
We can define correspondences more generally
replacing the fundamental class of $Z$ 
by any $\alpha \in H_{2d}^{BM}(Z)$.
In fact, the operators in our representation
shall be of this more general type.

\item If $Z$ is the graph of a proper morphism
$f:X_2 \to X_1$, then $c_{[Z]}=f_*$.
This shows that correspondences generalize
the pushforward in BM homology.

\item Since $X_1$ and $X_2$ are smooth,
we may use the isomorphism of BM and relative singular cohomology
to identify a class in the tensor product with an endomorphism
as we do in general for $V,W$ vector spaces via the usual isomorphism
$V \otimes W^* \cong \Hom(W,V)$, $v \otimes \varphi \mapsto 
(w \mapsto \varphi(w)v)$.
In this way the fundamental class 
$[Z] \in \bigoplus_{i+j=\dim_\mathbb{R}Z}H_i(X_1)\otimes H_j^{BM}(X_2)$
immediately gives a map
$$
\bigoplus_{n}H_* (X^{[n]}) \to \bigoplus_{n}H_* (X^{[n\mp i]})
$$
which, in fact, coincides with Equation \ref{equation-correspondence}
(see \cite[Section 8.3]{Nakajima-lectures}).

\end{enumerate}
\end{remark}

\medskip\noindent
Now we show how to compose correspondences,
as we shall do in the proof of the main theorem.

Suppose we are given two subvarieties
$Z_{12} \subset X_1 \times X_2$ and $Z_{23} \subset X_2 \times X_3$
such that the projections to $X_2$ and $X_3$ are proper.
We define the composition of the corresponding correspondences by
\begin{equation}
\label{equation-composition}
[Z_{12}]\circ[Z_{23}]:=p_{13,*}(p_{12}^* [Z_{12}]\cap p_{23}^* [Z_{23}])
\end{equation}
according to the diagram
$$
\xymatrix{
&  X_1 \times X_2 \times X_3\ar[dl]_{p_{12}}\ar[dr]^{p_{23}}\ar[dd]_{p_{13}}\\
 X_1\times X_2& & X_2\times X_3\\
& X_1\times X_3
}
$$
which is well-defined since, in fact,
the projection $p_{13}(p_{12}^{-1}Z_{12} \cap p_{23}^{-1} Z_{23})\to X_3$ 
is proper (see \cite[Section 8.3]{Nakajima-lectures}).




\section{Main construction}
\label{section-main-construction}

\noindent
In this section we define the
correspondences used in Nakajima's representation
of the Heisenberg algebra and formulate
the theorem which guarantees that it is so.

Let $n\geq i>0$. Define $P[i]$
as a subset of
$\coprod_nX^{[n-i]}\times X^{[n]}$
by
$$
P[i]:=\coprod_n \{(\mathcal{J}_1,\mathcal{J}_2,x)
\in X^{[n-i]}\times X^{[n]}:
\mathcal{J}_1 \supseteq \mathcal{J}_2,
\text{supp}(\mathcal{J}_1/\mathcal{J}_2)=\{x\}
\text{ for some }x \in X\}.
$$

\noindent
For $i<0$, we swap $\mathcal{J}_1$ and $\mathcal{J}_2$ 
in this definition.
Now let
$\alpha \in H_*^{BM}(X)$, $\beta \in H_*(X)$
and suppose that $i>0$. Define the classes
\begin{equation}
\label{equation-correspondence-p}
\begin{aligned}
P_\alpha[i]&=\varpi_*(\Pi_1^*\alpha \cap [P[i]])\\
P_\beta[i]&=\varpi_*(\Pi_1^*\beta \cap [P[-i]])
\end{aligned}
\end{equation}
according to the projections
$$
\xymatrix{
&  \coprod_n X^{[n-i]}\times X^{[n]}\times X
\ar[dl]_{\Pi_1}\ar[dr]^{\varpi_i}\\
X&&\coprod_n X^{[n-i]}\times X^{[n]}.
}
$$
These homology classes define the operators
of the group
$\bigoplus_{n=0}^{\infty}H_*(X^{[n]})$
that will be used in the main theorem.

\medskip\noindent
Before we continue, we formulate and
prove (with the help of other results,
as explained in \cite[Section 8.3]{Nakajima-lectures})
a lemma that will be used later.

\begin{lemma}
\label{lemma-dimension}
$\dim_\mathbb{C}P[i]\cap (X^{[n-i]}\times X^{[n]})=2n-i+1$.
\end{lemma}

\begin{proof}
Suppose $i>0$.
Fix $\mathcal{J}_1 \in X^{[n-i]}$.
By Theorem \ref{theorem-fogarty} we know that
$\dim X^{[n-i]}=2(n-i)$.
To find the dimension of $P[i]$
we must add to this number the dimension of the variety
obtained by varying  $\mathcal{J}_2 \in X^{[n]}$ 
satisfying $\mathcal{J}_1/\mathcal{J}_2=\{x\}$ 
for some $x \in X$.
The generic case is when 
$x \in X\setminus\text{supp}(\mathcal{O}_X/\mathcal{J}_1)$
(this may be accepted under the explanation
that since $\mathcal{O}_X/\mathcal{J}_1$ is the structure sheaf
of a subvariety, its support has less dimension than that of $X$;
a formal proof is provided in \cite[Lemma 8.31]{Nakajima-lectures}).
In fact, the set of such $\mathcal{J}_2$ is
$\pi^{-1}(S^i_{(i)}X)$, where $\pi$ is
the Hilbert-Chow morphism and
$S^i_{(i)}(X)=\{i[x] \in S^i(X):x \in X\}$.
This notation is used in other parts of \cite{Nakajima-lectures};
in our case just notice that $S^i_{(i)}(X)$
is the subset of the symmetric product $S^i(X)$ 
consisting of cycles of a single point with weight $i$.
Finally, $\pi^{-1}(S^i_{(i)}(X))$ has complex dimension $i+1$
as a subvariety of the Hilbert scheme.
\end{proof}

\medskip\noindent
Now we state the main theorem,
\cite[Theorem 8.13]{Nakajima-lectures}.
Notice that we must add the assumption that
the homology group of $X$ is concentrated in even degrees.
The theorem can be stated in a more general form without
this assumption, but we shall not discuss it in this work.
\begin{theorem}[Nakajima, Grojnowski]
\label{theorem-main}
Suppose $(\text{deg} \alpha)(\text{deg}\beta)\equiv 0 (\text{mod} 2)$
for all $\alpha,\beta \in H_*(X)$.
Then
\begin{equation}
\label{equation-main}
[P_\alpha[i],P_\beta[j]]=(-1)^{i-1}i \delta_{i,-j}
\left<\alpha,\beta\right>\text{Id}.
\end{equation}

\noindent
where $\left<\alpha,\beta\right>=p_*(\alpha \cap \beta)$
for $p:X \to \text{pt}$ the unique morphism
from $X$ to the point pt.
\end{theorem}

\medskip\noindent
According to our constructions in 
Section \ref{section-heisenberg},
this means that we have a representation of 
the Heisenberg algebra $\hat{\mathfrak{h}}$
associated to the graded vector space $H_*(X)$
on $\bigoplus_{n=0}^{\infty}H_{*}(X^{[n]})$.

\medskip\noindent
Notice the factor $(-1)^{i-1}$
in Equation \ref{equation-main},
which is not present in the Lie bracket equations of any of the
variants of the Heisenberg algebra discussed earlier in this work.
The reason for including this term in the main theorem 
is explained in
\cite[Theorem 9.15]{Nakajima-lectures}.
Moreover, this factor does not modify the Lie algebra structure:
we can define a Lie algebra isomorphism
from, say, the most basic version of the Heisenberg algebra
with bracket given by Equation \ref{equation-bracket-basic},
to the Lie algebra 
$\bigoplus_{n \in \mathbb{Z}}\mathbb{C}\tilde{h}_n \oplus\mathbb{C}K$ 
with bracket
$$
[\tilde{h}_m,\tilde{h}_m]=(-1)^{m-1}m\delta_{m,-n}K,\qquad 
\text{$K$ central,}
$$
by putting
$$
h_m \mapsto 
\begin{cases}
(-1)^{m-1}\tilde{h}_m\qquad &m<0 \\
\tilde{h}_m\qquad &m\geq 0.
\end{cases}
$$



\section{Worked example}
\label{section-example}

\noindent
In this section we do a worked example following \cite{Nakajima-lectures}
and \cite[Video 10]{videos}
to illustrate the main idea of the proof.

\medskip\noindent
Fix $x \in X$ and let $\alpha=[X]$ be the BM class of $X$,
and $\beta=[x]$ the class of the point $x$.
Take a set of distinct $n$ points
$x_1,\ldots,x_n \in X\setminus\{x\}$ 
none of which is $x$.
Identify them with the corresponding
ideal $\mathcal{J}$. We refer interchangeably
to the set of points and the ideal sheaf $\mathcal{J}$.

Let's show that
$$
[P_\alpha[1],P_\beta[-1]][\mathcal{J}]=[\mathcal{J}].
$$
We start by computing $P_\beta[-1][\mathcal{J}]$.
For this consider
$$
\xymatrix{
&X^{[n+1]}\times X^{[n]}\ar[ld]_{p_1}\ar[rd]^{p_2}\\
X^{[n+1]}&&X^{[n]}.
}
$$
By definition, $P_{[x]}[-1]$ is given by:
pull back $\mathcal{J}$ via $p_2$,
intersect with the class $\varpi_*(\Pi^*\alpha \cap [P[-1]])=P_{[x]}[-1]$ 
and push forward
down to $X^{[n+1]}$. 
This means that we push forward the cycle corresponding to
$$
\{(\mathcal{J}, \mathcal{J}',x)
:\mathcal{J}\subseteq \mathcal{J}',\text{supp}(\mathcal{J}'/\mathcal{J})=\{x\}\}
\subset X^{[n]} \times X^{[n+1]} \times X.
$$
The pushforward recovers the $X^{[n+1]}$ component,
that is, the schemes $\mathcal{J}'$ in $X^{[n+1]}$ satisfying
$\text{supp}(\mathcal{J}'/\mathcal{J})=\{x\}$.
Since $x$ is not a point in $\mathcal{J}$,
the only possible choice is that $\mathcal{J}'=\{x\}\cup \mathcal{J}$,
that is,
$$
P_{[x]}[-1][\mathcal{J}]=[\{x\}\cup \mathcal{J}].
$$
(see figure below).

On the other hand, the correspondence $P_{[X]}[1]$ 
starts with an ideal $\mathcal{J}'\in X^{[n+1]}$
intersects it with 
$$
\{(\mathcal{J}, \mathcal{J}',x)
:\mathcal{J}\subseteq \mathcal{J}',\text{supp}(\mathcal{J}'/\mathcal{J})=\{y\},
\text{ for some }y\in X\}
$$
and projects to $X^{[n]}$. We obtain
$\sum_{y \in \mathcal{J}'}[\mathcal{J}'-\{y\}]$.
Then it becomes clear that
$$
[P_\alpha[1],P_\beta[-1]][\mathcal{J}]=[\mathcal{J}].
$$
The following figure is taken directly from \cite{Nakajima-lectures}:

\includegraphics[width=0.9\textwidth]{fig1.png}
\label{figure-8.1}



\section{Proof of one case of the main theorem}
\label{section-proof}

\noindent
The proof of Theorem \ref{theorem-main} is split
in two different cases:
\begin{enumerate}
\item $i,j>0$ or $i,j<0$, and
\label{item-same}
\item $i>0$, $j<0$ or $i<0$, $j>0$.
\label{item-different}
\end{enumerate}

\noindent
In this work we explain the proof of case (\ref{item-same}).

\medskip\noindent
Suppose $i,j>0$ (the case $i,j<0$ is analogous).
Recalling our general definition for composition of correspondences
(Equation \ref{equation-composition})
we see that
the composition $P_\alpha[i]P_\beta[j]$ is given by
\begin{equation}
\label{equation-composition-pij}
P_\alpha[i]P_\beta[j]
=p_{13,*}(p_{12}^*[P[i]]\cap\Pi_1^*\alpha\cap p_{23}^*[P[j]]\cap\Pi_2^*\beta).
\end{equation}

Consider the set
\begin{equation}
\label{equation-cycle-pij}
\begin{aligned}
P[i,j]&:=p_{12}^{-1}(P[i])\cap p_{23}^{-1}(P[j])
\cap(X^{[n-i-j]}\times X^{[n-j]}\times X^{[n]})\\
&=\left\{(\mathcal{J}_1,\mathcal{J}_2,\mathcal{J}_3):
\mathcal{J}_1\supset\mathcal{J}_2\supset\mathcal{J}_3,
\substack{
\text{supp}(\mathcal{J}_1/\mathcal{J}_2)=\{x\}
\text{ for some }x\in X\\
\text{supp}(\mathcal{J}_2/\mathcal{J}_3)=\{y\}
\text{ for some }y \in X }\right\}
\end{aligned}
\end{equation}

\noindent
which corresponds to the diagram
$$
\xymatrix{
&  X^{[n-i-j]} \times X^{[n-j]} \times X^{[n]}
\ar[dl]_{p_{12}}\ar[dr]^{p_{23}}\ar[dd]_{p_{13}}\\
P[j]\subset X^{[n-j]}\times X^{[n]}& & X^{[n-i-j]}\times X^{[n-j]}\supset P[i]\\
& X^{[n-i-j]}\times X^{[n]}.
}
$$
By definition, the correspondence $P[j]$ maps the homology of $X^{[n]}$ to the 
homology of $X^{[n-j]}$. In order to compose this with $P[i]$
we must regard $P[i]$ as a map from the homology of $X^{[n-j]}$ 
to the homology of $X^{[n-j-i]}$, which is perfectly
valid since our correspondences are defined for all $n$.
Observe that the middle factor $X^{[n-j]}$
in the triple product will become $X^{[n-i]}$ when
we compute the composition in the opposite order.

\medskip\noindent
To recover the correspondence as in 
Equation \ref{equation-composition-pij}
from the set $P[i,j]$ it's natural to consider its 
fundamental class.
We will do this in the following special way.
Observe that $P[i,j]$ may be decomposed as a disjoint union
of the set $\{x \neq y\}$ 
of ``generic'' cases in which $x$ and $y$ in Equation \ref{equation-cycle-pij}
are different, and the set $\{x=y\}$ in which they coincide.
Let $\overline{P}[i,j]:=\overline{P[i,j]\cap\{x\neq y\}}$.
Then
$$
p_{12}^* [P[i]]\cap p_{23}^* [P[j]]=[\overline{P}[i,j]]+\iota_*R
$$
where $R \in H_*^{BM}(P[i,j]\cap \{x=y\})$ corresponds
to the $\{x=y\}$ part of $P[i,j]$ 
and $\iota$ is the inclusion.

\medskip\noindent
We can carry out the same construction
for the composition $P_\beta[j]P_\alpha[i]$,
obtaining
\begin{equation}
\label{equation-composition-pji}
P_\alpha[j]P_\beta[i]
=p_{13,*}(p_{12}^*[P[j]]\cap\Pi_2^*\alpha\cap p_{23}^*[P[i]]\cap\Pi_1^*\beta).
\end{equation}

\noindent
and
\begin{equation}
\label{equation-pji-r}
p_{12}^* [P[j]]\cap p_{23}^* [P[i]]=[\overline{P}[j,i]]+\iota'_*R'.
\end{equation}

\noindent
Observe that Equations \ref{equation-composition-pij} and
\ref{equation-composition-pji} differ in the classes
$\Pi_1^*\alpha$ against $\Pi_2^*\alpha$, and correspondingly with $\beta$.
The following lemma is all we need to obtain the result:

\begin{lemma}
\label{lemma-pij}
\begin{enumerate}
\item There exists an isomorphism
\begin{equation}
\label{equation-pij-isomorphism}
P[i,j]\cap\{x \neq y\}\xrightarrow{\cong}P[j,i]\cap \{x \neq y\},
\end{equation}

\noindent
which exchanges $\Pi_1$ and $\Pi_2$.

\item We have
\begin{equation}
\label{equation-r}
p_{13,*}(\Pi_1^*\alpha\cap\Pi_2^*\beta\cap\iota_*R)=0,\qquad 
p_{13,*}(\Pi_1\beta\cap\Pi_2^*\alpha\cap\iota_*R')=0.
\end{equation}
\end{enumerate}
\end{lemma}

\noindent
Indeed, this lemma is all we need since
\begin{align*}
P_\alpha[i]P_\beta[j]
&=(-1)^{\text{deg}\alpha}p_{13,*}((\overline{P}[i,j]+\iota_*R)\cap\Pi_1^*\alpha\cap\Pi_2^*\beta)\\
&=(-1)^{\text{deg}\alpha\text{deg}\beta}
p_{13,*}((\overline{P}[j,i]+\iota_*R')\cap\Pi_1^*\beta\cap\Pi_2^*\alpha)\\
&=P_\beta[j]P_\alpha[i].
\end{align*}

\begin{proof}[Proof of Lemma \ref{lemma-pij}]
\begin{enumerate}
\item In this proof I assume that ``isomorphism'' means bijection, since
$P[i,j]$ is considered as a set-theoretical intersection.

First notice that when we say that the isomorphism, say $\varphi$, 
``exchanges $\Pi_1$ and $\Pi_2$''
we mean that $\Pi_1|_{X \times X}=\Pi_2|_{X \times X}\circ\varphi$
according to the diagram
$$
\xymatrix{
& & X\\
P[i,j] \cap \{x\neq y\}\subset
\ar[ddd]_\varphi&
\save{
X^{[n-i-j]}\times X^{[n-j]} \times X^{[n]} \times X \times X
}
\ar[ur]^{\Pi_1}\ar[dr]_{\Pi_2}
\restore
\\
& & X\\
& & X\\
P[j,i]\cap\{x\neq y\}\subset&
X^{[n-i-j]}\times X^{[n-i]} \times X^{[n]} \times X \times X
\ar[ur]^{\Pi_1}\ar[dr]_{\Pi_2}\\
& & X.
}
$$

Now, the point of the isomorphism $\varphi$ is that it must
map an ideal in $X^{[n-j]}$ to an ideal in $X^{[n-i]}$.
More precisely, we map 
$(\mathcal{J}_1,\mathcal{J}_2,\mathcal{J}_3)\in P[i,j] \cap \{x \neq  y\}$
to $(\mathcal{J}_1,\mathcal{J}_2',\mathcal{J}_3) \in P[j,i]\cap \{x \neq y\}$.
The defining relation is
\begin{equation}
\label{equation-isomorphism}
\mathcal{J}_1/\mathcal{J}_2'=\mathcal{J}_2/\mathcal{J}_3,\qquad 
\mathcal{J}_2'/\mathcal{J}_3=\mathcal{J}_1/\mathcal{J}_2.
\end{equation}

\noindent
The following figure 
(taken from \cite{Nakajima-lectures})
shows how this correspondence works.

\includegraphics[width=0.8\textwidth]{fig2.png}

The setting is the following. We have the ideals with contentions
$\mathcal{J}_1\supset \mathcal{J}_2,\mathcal{J}_2'\supset \mathcal{J}_3$.
Passing to the supports, the inclusions are reversed, i.e 
the support of $\mathcal{J}_3$ contains the rest,
and the support of $\mathcal{J}_1$ is contained in the rest.

To prove that this function is well-defined
we need to show that for every $\mathcal{J}_2$ there is a unique
$\mathcal{J}_2'$ satisfying Equation \ref{equation-isomorphism}.
Notice that by definition, $\text{supp}\mathcal{J}_2/\mathcal{J}_3=\{x\}$
for some $x \in X$ so that, as a scheme, $\{x\}$ has multiplicity $j$.
We define
$\mathcal{J}_2'=\mathcal{J}_1\cup \{x\}$
where $x$ has multiplicity $j$.

\medskip\noindent
(Caveat: it looks to me that the choice of $\mathcal{J}_2'$ is not unique.
Indeed, the space of schemes of length $j$ supported
in the point $x$ was discussed in the proof of Lemma \ref{lemma-dimension},
where we used the fact that it has dimension $j-1$.)

\medskip\noindent
Defining $\mathcal{J}'_2$ this way implies that
$$
\mathcal{J}_1/\mathcal{J}_2'=\mathcal{J}_1/(\mathcal{J}_1\cup \{x\})
=\{x\}=\mathcal{J}_2/\mathcal{J}_3.
$$
Now, we also must have $\mathcal{J}_1/\mathcal{J}_2=\{y\}$ 
for some $y \in X$ with multiplicity $i$.
This implies that $\mathcal{J}_1/\mathcal{J}_3=\{x\}\cup \{y\}$ 
with their multiplicities.
Then 
$$
\mathcal{J}_2'/\mathcal{J}_3=(\mathcal{J}_1\cup \{x\})/\mathcal{J}_3=\{y\}
=\mathcal{J}_1/\mathcal{J}_2
$$
as required.
An inverse function 
defined similarly shows this is a bijection.

\item Consider the commutative diagram
$$
\xymatrix{
P[i,j]\cap \{x=y\}\ar[r]^-{\Pi'}\ar[d]_{\iota}&X\ar[d]^{\Delta}\\
P[i,j]\ar[r]_{\Pi_1 \times \Pi_2}&X \times X
}
$$
where $\Delta$ is the diagonal embedding
and $\Pi'$ is any of the projections $\Pi_1$ or $\Pi_2$
(which coincide on  $P[i,j]\cap \{x=y\}$).
Then
\begin{equation}
\label{equation-lemma2-1}
\begin{aligned}
\Pi_1^*\alpha \cap \Pi_2^*\beta \cap\iota_*R
&=\iota_*(\iota^*(\Pi_1^*\alpha\cap\Pi_2^*\beta)\cap R)\\
&=\iota_*(\iota^*(\Pi_1\times\Pi_2)^*(\alpha\times\beta)\cap R)\\
&=\iota_*({\Pi'}^*\Delta^*(\alpha\times\beta)\cap R)\\
&=\iota_*({\Pi'}^*(\alpha\cap\beta)\cap R),
\end{aligned}
\end{equation}

\noindent
where the first equality is due
to the projection formula
(see Stacks Project 
\href{https://stacks.math.columbia.edu/tag/01E8}{01E8})
and we stick with Nakajima's notation $\alpha \times \beta$ 
for the tensor product $\alpha \otimes \beta$.
We also have the commutative diagram
$$
\xymatrix{
p_{13}\circ \iota:P[i,j]\cap\{x=y\}
\ar[r]^-{\Pi'}\ar[d]_{p_{13 \circ \iota}}&X\ar@{=}[d]\\
P[i+j]=\{(\mathcal{J}_1,\mathcal{J}_3)|\mathcal{J}_1 \supset \mathcal{J}_3,
\text{supp}(\mathcal{J}_1/\mathcal{J}_3)=\{x\}\}\ar[r]_-{\Pi''}&X,
}
$$
where $\Pi''$ is the usual projection to $X$ from the
set $P[i+j]$.

By Equation \ref{equation-lemma2-1}
and the commutativity of the second diagram,
we can express the class of interest as intersection
of two classes: one pulled back directly from $X$,
and the other one corresponding to $R$.
That is, we have
$$
p_{13,*}\iota_*({\Pi'}^* (\alpha\cap \beta)\cap R)
={\Pi''}^* (\alpha \cap \beta)\cap p_{13,*}\iota_*(R).
$$
Now, by Lemma \ref{lemma-dimension},
we know that
$$
\dim_{\mathbb{C}}P[i+j]=2n-i-j+1.
$$
On the other hand, the ``expected dimension''
is $2n-i-j+2$. This obliges $\iota_*R$ to vanish.

\medskip\noindent
(Caveat: I understand that the ``expected dimension'' 
means in \cite{Nakajima-lectures} the dimension of
$P[i,j]\cap \{x\neq y\}$, which is claimed to be
$2n-i-j+2$.
It is however not clear to me why this would imply
that $p_{13,*}\iota_*R$ must vanish.)

\end{enumerate}
\end{proof}

\bibliography{my}
\bibliographystyle{amsalpha}

\end{document}
